\section{Les IAG et la remise en question des \'evaluations}\label{sec:remise}

\subsection{Transformation des pratiques d'\'evaluation}
% - Devoirs, dissertations, exercices
% - Problème de l'authenticité

Les modalit\'es d'\'evaluation traditionnelles --- dissertation, compte rendu, exercice \`a domicile, m\'emoire --- ont \'et\'e con\c{c}ues dans un environnement o\`u la production \'ecrite personnelle constituait une t\^ache forc\'ement individuelle, n\'ecessitant un effort cognitif soutenu: recherche documentaire, organisation des id\'ees, ma\^\i trise de la langue, argumentation. 
Cet effort \'etait lui-m\^eme formateur: le processus de r\'edaction contribuait \`a la construction des comp\'etences que l'on cherchait \`a \'evaluer.

Les IAG transforment cet environnement. 
Il est d\'esormais possible de g\'en\'erer en quelques minutes un texte structur\'e, argument\'e et bien r\'edig\'e sur presque n'importe quel sujet. 
C'est le probl\`eme de l'\emph{authenticit\'e}: comment savoir si un travail remis est le produit de l'effort cognitif de l'apprenant, ou le r\'esultat d'une d\'el\'egation \`a une machine~?

Le cas du m\'emoire est particuli\`erement significatif, car c'est l'exercice qui condense le plus d'exigences intellectuelles --- recherche, synth\`ese, argumentation originale. 
Comme le montrent les \'etudes men\'ees par Leroux et al.\ sur les m\'ethodes d'\'evaluation au coll\'egial, les instruments d'\'evaluation ont \'et\'e con\c{c}us pour rendre compte d'un \emph{parcours de d\'eveloppement} individuel~\citep{leroux2015evaluer}.
Or, l'IAG vient perturber ce parcours en s'ins\'erant comme interm\'ediaire entre l'\'etudiant et sa production. 
Une recherche men\'ee \`a l'intersection des travaux sur les IAG et les m\'emoires universitaires soul\`eve la question suivante: lorsque l'\'etudiant utilise l'IAG pour r\'ediger une partie de son travail, en quoi le m\'emoire t\'emoigne-t-il encore de \emph{ses} comp\'etences propres?

La r\'eponse des institutions a souvent consist\'e \`a tenter de d\'etecter les productions IA via des logiciels anti-plagiat. 
Mais cette approche est limit\'ee: les d\'etecteurs produisent des taux importants de faux positifs et de faux n\'egatifs, et leur fiabilit\'e est insuffisante pour fonder des d\'ecisions acad\'emiques. 
De plus, ils cr\'eent une course aux armements entre les outils de d\'etection et les outils de g\'en\'eration, sans traiter la question de fond.

\subsection{D\'eplacement de la notion de comp\'etence}
% - compétence cognitive vs compétence assistée
% - rôle de l'IA dans la production de savoir

L'usage des IAG soul\`eve une question philosophique sur la nature m\^eme de la comp\'etence: peut-on parler de \emph{comp\'etence assist\'ee} au m\^eme titre que de comp\'etence individuelle~?

Ce d\'ebat n'est pas sans pr\'ec\'edent. 
L'introduction de la calculatrice dans les classes de math\'ematiques avait provoqu\'e des tensions analogues: calculer \`a la main et calculer avec une machine mobilisent des comp\'etences diff\'erentes. 
La calculatrice a finalement \'et\'e int\'egr\'ee, avec un recentrage de l'\'evaluation sur la mod\'elisation et l'interpr\'etation plut\^ot que sur l'op\'eration arithm\'etique. 
Avec les IAG, le d\'efi est d'une autre ampleur: leur p\'erim\`etre d'action chevauche les comp\'etences d'ordre sup\'erieur que l'\'ecole vise \`a d\'evelopper --- l'argumentation, la synth\`ese, le raisonnement critique.
Il est alors difficile de déterminer vers quoi recentrer l'\'evaluation si l'outil couvre déjà une grande partie du domaine traditionnellement évalué.

Les travaux sur les comp\'etences \`a acqu\'erir mettent en avant la n\'ecessit\'e de \emph{m\'etacomp\'etences}: apprendre \`a apprendre, savoir s'informer, \'evaluer la pertinence et la fiabilit\'e des sources~\citep{heiser2023education}. 
Or, une utilisation non critique des IAG risque pr\'ecis\'ement de court-circuiter le d\'eveloppement de ces m\'etacomp\'etences: l'\'etudiant qui recopie une synth\`ese IAG sans l'\'evaluer n'apprend pas \`a \'evaluer. 
\`A l'inverse, l'\'etudiant qui utilise l'IAG comme outil dialectique --- pour tester des arguments, g\'en\'erer des contre-exemples, structurer sa pens\'ee --- peut exercer des comp\'etences \'elabor\'ees.

Selon les travaux sur l'impact des IAG sur le d\'eveloppement des comp\'etences, la distinction cruciale est celle entre deux r\^oles que l'IA peut jouer: \emph{substitut} (elle remplace le processus cognitif de l'apprenant) ou \emph{outil} (elle amplifie et soutient ce processus)~\citep{shiohira2021comprendre}. 
Cette distinction, claire en th\'eorie, est difficile \`a op\'erationnaliser dans la pratique \'evaluative: le rendu final ne permet g\'en\'eralement pas de d\'eterminer lequel de ces deux r\^oles l'IA a jou\'e.
